\partstart{Thảo luận và so sánh}
\subsection{K-NN và cây quyết định trên dữ liệu hỗn hợp}
K-NN biểu diễn giả thuyết rằng các đơn gần nhau trong không gian đầu vào
có xu hướng chung nhãn. Điểm mạnh là không cần định sẵn dạng ranh giới
tuyến tính. Điểm yếu trong bài toán này là sự gần nhau phụ thuộc mạnh vào
mã hóa và thang đo. Một không gian nhiều cột one-hot có thể làm khoảng
cách khó diễn giải về nghiệp vụ; các trường địa lý hiếm còn ít láng giềng
để tham chiếu đáng tin cậy.

Cây quyết định tìm những điều kiện tách dữ liệu thành nhóm tương đối đồng
nhất. So với K-NN, đường đi qua các nút dễ đọc hơn và chi phí dự đoán trên
cây nông ít phụ thuộc số dòng train. Tuy vậy, quy tắc dễ đọc không đồng
nghĩa đúng nhân quả hay ổn định. Một split nhỏ có thể thay đổi nhánh; cây
quá sâu có thể khớp những nhóm hiếm và thể hiện khoảng cách train--CV lớn
như baseline đã quan sát.

Trong thí nghiệm, cây B vượt K-NN B về Precision và Accuracy nhưng Recall
thấp hơn. Không có thuật toán thắng toàn bộ tiêu chí. Nếu mục tiêu là
ưu tiên một danh sách ngắn có tỷ lệ đúng cao, hướng đánh đổi của cây B
đáng khảo sát thêm; nếu mục tiêu là giảm bỏ sót, Recall thấp hiện tại là
hạn chế rõ. Chưa có yêu cầu nghiệp vụ định lượng để lựa chọn giữa hai
hướng, nên báo cáo chỉ giữ tiêu chí CV F1 đã định trước.

\subsection{Ảnh hưởng quan sát được của luồng A và B}
Bổ sung lịch có lợi về F1 trong các so sánh đã lưu, nhưng mức cải thiện
khác nhau theo thuật toán và cách chọn tham số. K-NN có chênh lệch CV
0,0046, còn cây tuned có chênh lệch 0,0316. Khi giữ tham số cây A cố định,
chênh lệch chỉ 0,0163. Vì vậy, không quy toàn bộ mức tăng của cây tuned
cho ba cột thời gian; một phần so sánh đồng thời thay đổi độ sâu được chọn.

Tháng, thứ, giờ có thể là proxy cho chu kỳ vận hành, cơ cấu dịch vụ hoặc
giai đoạn thu thập dữ liệu. Đây là các khả năng giải thích cần kiểm chứng,
không phải nguyên nhân đã xác nhận. Chẳng hạn, giờ đặt đơn quan trọng trong
cây có thể phản ánh cơ chế dữ liệu riêng của Same Day thay vì một quy luật
chung cho mọi dịch vụ. Muốn kết luận cần thêm dữ liệu và thí nghiệm được
định trước, đặc biệt giữ tập tương lai độc lập.

Hai luồng không phải so sánh ``có tiền xử lý'' với ``không tiền xử lý''.
Cả hai đều có imputation, scaling và one-hot. Do đó, kết quả chỉ hỗ trợ
đối chiếu biểu diễn A với A cộng lịch, không đo riêng hiệu quả từng phép
biến đổi. Đây là giới hạn thiết kế cần nêu để tránh mô tả phương pháp
nhiều hơn bằng chứng thực sự có.

\subsection{Baseline, Dummy và ý nghĩa của F1 cạnh tranh}
Với tỷ lệ dương $\pi$ và quy tắc luôn dự đoán dương, Precision bằng $\pi$,
Recall bằng 1, nên
\begin{equation}
F_{1,\mathrm{all+}}=\frac{2\pi}{1+\pi}.
\end{equation}
Thay $\pi=632/1176$ cho test dòng thu được F1 khoảng 0,6991, đúng với
Dummy. Đây là hệ quả toán học từ số đếm lớp, không cần mô hình khai thác
đặc trưng. Vì vậy, một điểm F1 xấp xỉ 0,70 chưa tự chứng minh mô hình
phân biệt hữu ích trong bài toán này.

Cây B đạt F1 dòng cao hơn Dummy khoảng 0,0107 nhưng Precision tăng từ
0,5374 lên 0,8533, còn Recall giảm từ 1 xuống 0,6076. Đó là thay đổi
thực chất về kiểu ra quyết định. Cây chỉ cảnh báo 450/1.176 dòng, trong
khi Dummy cảnh báo toàn bộ. Tỷ lệ cảnh báo đúng cao hơn có thể giúp hạn
chế khối lượng xử lý, nhưng 248 FN vẫn là phần rủi ro lớn cần cân nhắc.

Cây baseline cho thấy một kiểu thất bại khác: gần hoàn hảo trên train
nhưng không vượt Dummy ở CV hoặc test F1. Tuned tree A cải thiện ít so
với baseline về CV và vẫn dưới Dummy. Như vậy, sử dụng thuật toán có học
không tự bảo đảm vượt đối chứng không dùng đặc trưng. Việc báo đủ baseline
là cần thiết để không trình bày metric rời bối cảnh.

\subsection{False Positive và False Negative trong nghiệp vụ}
FP là cảnh báo cho dòng được nguồn gán âm. Nó có thể dẫn tới cuộc gọi,
kiểm tra kho hoặc ưu tiên không cần thiết, nhưng dữ liệu không ghi nhận
biện pháp nào đã được thực hiện. Hơn nữa, nhãn âm còn có đơn hủy, nên
không mô tả mọi FP là ``đơn đã giao đúng hạn bị báo nhầm''. Đó là lý do
ngữ nghĩa nhãn ảnh hưởng trực tiếp tới phân tích tác động.

FN là dòng trễ không được cảnh báo. Với cây B, 248 FN trên 632 dòng dương
tương đương bỏ sót khoảng 39,24\% dòng dương toàn test. Người vận hành
không nhìn thấy các dòng này trong danh sách cảnh báo. Tuy nhiên, không
biết một can thiệp sớm có thực sự ngăn trễ hay không, nên không chuyển
số FN thành số đơn chắc chắn có thể cứu được hoặc thiệt hại tiền cụ thể.

Nếu ký hiệu chi phí một FP là $c_{FP}$ và một FN là $c_{FN}$, một mô hình
đơn giản cho tổng lỗi là $C=c_{FP}FP+c_{FN}FN$. So cây B với Dummy trên
test dòng cho
\begin{equation}
C_{B}-C_{D}=-478c_{FP}+248c_{FN}.
\end{equation}
Trong mô hình chi phí giả định này, B có chi phí thấp hơn khi
$c_{FN}/c_{FP}<478/248\approx1{,}927$, nếu $c_{FP}>0$. Đây là phép suy
ra minh họa từ số đếm đã có, không phải kết quả kinh tế được đo. Nó bỏ qua
chi phí xử lý TP, khả năng can thiệp thành công và phụ thuộc giữa các dòng.

Do nhiều dòng thuộc cùng đơn, chi phí theo dòng có thể tính lặp một hành
động. Nếu quyết định ở cấp đơn, cần dùng số FP/FN theo đơn và chi phí phù
hợp cấp đơn. Không chọn tỷ lệ chi phí sau khi thấy kết quả để chứng minh
mô hình đã tốt. Chi phí và mục tiêu Recall nên được thống nhất trước
thiết kế validation và lựa chọn ngưỡng tiếp theo.

\subsection{Standard Class là điểm yếu trọng yếu}
Standard Class chiếm 704/1.176 dòng test và chứa 254 dòng dương. Cây B
chỉ phát hiện sáu dòng dương của nhóm, bỏ sót 248. Điều này tạo nghịch lý
bề ngoài: Precision toàn mô hình cao nhưng một dịch vụ chiếm phần lớn
test gần như không được phát hiện giao trễ. Metric tổng hợp bị chi phối
bởi việc cây cảnh báo tốt ở các dịch vụ khác.

Tỷ lệ âm cao trong Standard Class khiến quy tắc phần lớn dự đoán âm có
thể đạt nhiều TN. Đây là giải thích phù hợp với số đếm và cấu trúc cây,
không khẳng định là nguyên nhân duy nhất của mọi FN. Nhãn của các đơn,
thời hạn kế hoạch và cách tạo dữ liệu có thể cùng ảnh hưởng. Không có
bằng chứng ở artifact để kết luận kho, tuyến đường hay nhân viên cụ thể
gây ra các trường hợp giao trễ này.

Ngược lại, First Class và Second Class có toàn bộ dòng âm bị cảnh báo
sai trong holdout. Hai nhóm có FPR bằng 1 nhưng mẫu số âm lần lượt chỉ
10 và 40. Báo cả tử số lẫn mẫu số giúp tránh đọc tỷ lệ như quy luật phổ
quát. Same Day không lỗi trên 71 dòng vẫn cần kiểm tra ở dữ liệu khác;
không coi số lỗi bằng 0 ở một mẫu là bằng chứng bảo đảm dịch vụ.

Một hướng nghiên cứu hợp lý là kiểm tra mô hình hoặc tiêu chí validation
có chú ý Recall từng dịch vụ. Đây là đề xuất tương lai: hiện chưa chạy
mô hình riêng Standard Class, chưa tối ưu trọng số lớp và chưa chọn
threshold theo nhóm. Những thay đổi đó cần được thiết kế, fit và lựa chọn
trên train/validation trước một test độc lập, tránh học lại từ test hiện tại.

\subsection{Diễn giải đặc trưng và quy tắc cây}
\begin{table}[!htbp]\centering
\caption{Bốn importance lớn nhất trong cây B theo artifact đã lưu.}\label{tab:importance}
\begin{tabularx}{\linewidth}{|Y|r|}\hline
\theadcell{Đặc trưng sau tiền xử lý} & \theadcell{Importance}\\\hline
Shipping Mode: Standard Class & 0,5735\\
Giờ tạo đơn & 0,1758\\
Shipping Mode: Same Day & 0,1356\\
Shipping Mode: First Class & 0,0624\\\hline
\end{tabularx}

\end{table}
Bảng~\ref{tab:importance} cho thấy phần lớn mức giảm impurity gắn với
Shipping Mode và giờ tạo đơn. Quy tắc xuất bắt đầu bằng chỉ báo Standard
Class, sau đó có các nhánh Same Day và địa lý. Điều này phù hợp với sự
khác nhau rõ về hành vi dự đoán giữa nhóm dịch vụ. Importance thuộc từng
cột one-hot, không phải phép đo nhân quả của toàn trường gốc.

Một phần quy tắc thực trong \code{tree_B_rules.txt} là nhánh Same Day
so sánh \code{numeric__order_hour <= 0.03}. Giá trị 0,03 nằm trên thang
StandardScaler, không phải 0,03 giờ trong ngày. Muốn chuyển về giờ gốc
cần dùng $x=\mu+s z$ với tham số scaler đã fit của đúng mô hình. Báo cáo
không có đủ các tham số đó trong bảng aggregate để tự thay ngưỡng bằng
một giờ cụ thể.

Các nhánh có thể chứa điều kiện khác nhau nhưng kết thúc cùng nhãn 1.
Điều đó không có nghĩa điều kiện hoàn toàn vô dụng: phân bố xác suất ở
lá có thể khác dù argmax giống nhau. Ngược lại, chỉ xem nhãn ở cây xuất
không đủ để khẳng định xác suất chính xác. Muốn dùng xác suất làm ưu tiên
cần kiểm tra calibration và độ ổn định ở dữ liệu độc lập.

\subsection{Đánh giá cấp dòng và cấp đơn}
Metric dòng trả lời hiệu quả trên một dòng hàng được chọn theo phân phối
test. Metric đơn trả lời hiệu quả khi mỗi đơn có trọng số ngang nhau.
Nếu đơn nhiều mặt hàng gặp cùng lỗi, metric dòng có thể khuếch đại trọng
số của tình huống đó. Chia nhóm giải quyết việc rò rỉ giữa hai tập, nhưng
không tự thay đổi trọng số metric từ dòng sang đơn.

Trên 423 đơn, cây B có Precision 0,8554, Recall 0,6068 và F1 0,7100;
Dummy F1 0,7123. Mức chênh nhỏ và thứ tự đảo so với F1 dòng làm rõ việc
chọn đơn vị đánh giá phải gắn với đơn vị hành động. Không thay tiêu chí
chọn mô hình sau khi nhìn thấy sự đảo thứ tự. Trong nghiên cứu tiếp theo,
nếu sản phẩm cảnh báo theo đơn thì metric validation cũng nên được xác
định ở cấp đơn ngay từ đầu.

Quy tắc trung bình xác suất giả định mọi dòng của đơn có vai trò ngang
nhau. Một mặt hàng đặc biệt có thể cần cách tổng hợp khác, nhưng dataset
không cung cấp mục tiêu hành động đủ rõ để chọn quy tắc mới. Không suy
ra trung bình tốt hơn max hoặc ngược lại khi chưa có thí nghiệm. Phân
tích hiện tại chỉ kiểm tra một cách tổng hợp đã được lưu trong protocol.

\subsection{Tính đại diện, giới hạn dữ liệu và tổng quát hóa}
Mẫu 2.113 đơn là phần nhỏ trong 65.752 đơn nguồn. Kiểm tra sampling ghi
nhận chênh lệch tỷ lệ biên lớn nhất khoảng 0,0227 ở Customer Segment
trong các danh mục đã khảo sát. Sự gần nhau của một số tỷ lệ không bảo
đảm quan hệ giữa phương thức giao, vùng, thời gian và nhãn cũng giống
nhau. Các tổ hợp hiếm có thể còn rất ít hoặc vắng trong mẫu.

Cách chọn tiền tố sau xáo trộn nhóm giữ nguyên đơn nhưng vẫn chịu ràng
buộc kích thước nhóm và ngân sách. Không gọi nó là mẫu phân tầng theo
mọi trường hoặc mẫu không chệch đã chứng minh. Một seed duy nhất không
đo độ nhạy kết quả với việc đơn nào được đưa vào. Lặp lấy mẫu theo nhóm
có thể cung cấp phân phối metric, nhưng chưa được thực hiện ở đây.

Nguồn lịch sử còn có khác biệt giữa nhãn giao trễ và trạng thái hủy,
giới hạn về cách ghi timestamp, danh mục địa lý và thời điểm một số trường
được hình thành. Không có kiểm chứng vận hành để xác nhận các quan hệ
quan sát được sẽ giữ nguyên ở doanh nghiệp khác. Vì vậy, không dùng từ
``tổng quát tốt'' chỉ dựa trên một holdout nhóm ngẫu nhiên.

\subsection{Khả năng ứng dụng và những hạn chế còn lại}
Ở trạng thái hiện tại, sản phẩm phù hợp làm nghiên cứu học phần và bằng
chứng cho quy trình có kiểm soát. Một ứng dụng thử nghiệm hợp lý trong
tương lai là chạy song song với quy trình hiện hữu, chỉ ghi cảnh báo để
đo chất lượng mà chưa tự điều phối. Cần thống nhất thời điểm chấm điểm,
đơn vị đơn hàng, định nghĩa trễ và mục tiêu phát hiện trước khi triển khai.

Các điều kiện còn thiếu gồm kiểm định thời gian, dữ liệu vận hành độc
lập, chi phí lỗi, hiệu chuẩn xác suất, khoảng tin cậy theo nhóm và giám
sát drift. Ngoài ra, nhóm cần xác nhận tính sẵn có của thuộc tính ngay
tại lúc tạo đơn. Một mô hình chạy đúng pipeline nhưng dùng cột xuất hiện
muộn hơn giả định vẫn không đáp ứng bài toán cảnh báo sớm.

Không có kết quả nào trong báo cáo chứng minh cảnh báo giúp giảm tỷ lệ
giao trễ thực tế. Để kiểm tra điều đó cần thiết kế đo tác động can thiệp,
khác với bài toán đo dự đoán. Việc giữ ranh giới giữa hiệu quả phân loại
và hiệu quả điều hành giúp các hướng phát triển có tiêu chí kiểm chứng
rõ ràng, thay vì chỉ đặt mục tiêu nâng thêm vài phần trăm F1.
